% GENERATED FILE. Edit canonical lessons, then run npm run generate:books.
% canonical-source-hash: fnv1a64:09205ea8e7537ddb
% canonical-lessons: SA-C550-vantayati, SA-C550-pratyekam, SA-C550-dayah, SA-C550-viniyunkte, SA-C550-vitaranam

\chapter{Sharing Out}
\label{ch:sa-sharing-out}
\hlchaptermodality{550}

\begin{chapteropening}
\textbf{By the end of this chapter:} \emph{I can say how things will be shared if everyone agrees.}

Complete the last lesson of chapter 550: I can say how things will be shared if everyone agrees.
\end{chapteropening}

\section[\texorpdfstring{\sk{वण्टयति} (vaṇṭayati)}{वण्टयति (vaṇṭayati)}]{\sk{वण्टयति} (vaṇṭayati) — divides up; shares out}

\label{lesson:SA-C550-vantayati}



\begin{quote}
Before the new one: say the Sanskrit for a wish to see, then the Sanskrit for strives for.
\end{quote}

\subsection*{You'll want to know: \sk{वण्टयति}}
\textbf{\sk{वण्टयति}} — \emph{vaṇṭayati} — \textquotedblleft{}divides up; shares out\textquotedblright{}.

\begin{quote}
\sk{यदि} \sk{माता} \sk{मोदकान्} \sk{वण्टयति}, \sk{तर्हि} \sk{सर्वे} \sk{एकं} \sk{लभन्ते।}
\end{quote}

\begin{quote}
\emph{yadi mātā modakān vaṇṭayati, tarhi sarve ekaṃ labhante.}
\end{quote}

\begin{quote}
If my mother shares out the sweets, everyone gets one.
\end{quote}

Past (she …): \sk{सा} \sk{मोदकान्} \sk{वण्टितवती}, \emph{sā modakān vaṇṭitavatī}.

\begin{cousinweb}
\textbf{\sk{वण्टयति}} (\emph{vaṇṭayati}) comes from the root √vaṇṭ, \textquotedblleft{}to partition, to share, to divide\textquotedblright{}; its past participle is \emph{vaṇṭita} (MW).
\end{cousinweb}

\begin{grammarlens}[title={What you've built}]
One more word for saying what happens if.
\end{grammarlens}

\subsection*{Your turn}
\begin{itemize}
\raggedright
  \item \emph{Say it:} \emph{vaṇṭayati}
  \item \emph{Say it:} \emph{vaṇṭayati}, once more
  \item \emph{Say it:} \emph{yadi mātā modakān vaṇṭayati, tarhi sarve ekaṃ labhante.} — If my mother shares out the sweets, everyone gets one.
  \item \emph{Recall:} say the Sanskrit for a wish to see, then the Sanskrit for strives for, then say \emph{vaṇṭayati} again
\end{itemize}

\begin{tcolorbox}[breakable,colback=teal!4,colframe=teal!35!black,title={Before you move on}]
What does \sk{वण्टयति} mean? (\textquotedblleft{}Divides up; shares out\textquotedblright{}.) Say it once more.
\end{tcolorbox}
\section[\texorpdfstring{\sk{प्रत्येकम्} (pratyekam)}{प्रत्येकम् (pratyekam)}]{\sk{प्रत्येकम्} (pratyekam) — each one; one by one}

\label{lesson:SA-C550-pratyekam}



\begin{quote}
Before the new one: say the Sanskrit for strives for, then the Sanskrit for divides up.
\end{quote}

\subsection*{You'll want to know: \sk{प्रत्येकम्}}
\textbf{\sk{प्रत्येकम्}} — \emph{pratyekam} — \textquotedblleft{}each one; one by one\textquotedblright{}.

\begin{quote}
\sk{यदि} \sk{दश} \sk{फलानि} \sk{सन्ति}, \sk{तर्हि} \sk{प्रत्येकं} \sk{बालः} \sk{एकं} \sk{लभते।}
\end{quote}

\begin{quote}
\emph{yadi daśa phalāni santi, tarhi pratyekaṃ bālaḥ ekaṃ labhate.}
\end{quote}

\begin{quote}
If there are ten fruits, each child gets one.
\end{quote}

\begin{cousinweb}
\textbf{\sk{प्रत्येकम्}} (\emph{pratyekam}) is built from \emph{prati}, \textquotedblleft{}towards, for each\textquotedblright{}, and \emph{eka}, \textquotedblleft{}one\textquotedblright{}; MW gives \emph{pratyekam} as one by one, singly, for every single one (MW).
\end{cousinweb}

\begin{grammarlens}[title={What you've built}]
One more word for saying what happens if.
\end{grammarlens}

\subsection*{Your turn}
\begin{itemize}
\raggedright
  \item \emph{Say it:} \emph{pratyekam}
  \item \emph{Say it:} \emph{pratyekam}, once more
  \item \emph{Say it:} \emph{yadi daśa phalāni santi, tarhi pratyekaṃ bālaḥ ekaṃ labhate.} — If there are ten fruits, each child gets one.
  \item \emph{Recall:} say the Sanskrit for strives for, then the Sanskrit for divides up, then say \emph{pratyekam} again
\end{itemize}

\begin{tcolorbox}[breakable,colback=teal!4,colframe=teal!35!black,title={Before you move on}]
What does \sk{प्रत्येकम्} mean? (\textquotedblleft{}Each one; one by one\textquotedblright{}.) Say it once more.
\end{tcolorbox}
\section[\texorpdfstring{\sk{दायः} (dāyaḥ)}{दायः (dāyaḥ)}]{\sk{दायः} (dāyaḥ) — an inheritance}

\label{lesson:SA-C550-dayah}



\begin{quote}
Before the new one: say the Sanskrit for divides up, then the Sanskrit for each one.
\end{quote}

\subsection*{You'll want to know: \sk{दायः}}
\textbf{\sk{दायः}} — \emph{dāyaḥ} — \textquotedblleft{}an inheritance\textquotedblright{}.

\begin{quote}
\sk{यदि} \sk{पितुः} \sk{दायः} \sk{अस्ति}, \sk{तर्हि} \sk{भ्रातरः} \sk{तं} \sk{वण्टयन्ति।}
\end{quote}

\begin{quote}
\emph{yadi pituḥ dāyaḥ asti, tarhi bhrātaraḥ taṃ vaṇṭayanti.}
\end{quote}

\begin{quote}
If there is an inheritance from their father, the brothers share it.
\end{quote}

A masculine noun.

\begin{cousinweb}
\textbf{\sk{दायः}} (\emph{dāyaḥ}) comes from the root √dā, \textquotedblleft{}to give\textquotedblright{}, which also gives \emph{dadāti}; MW gives \emph{dāya} as a share, a portion, an inheritance (MW).
\end{cousinweb}

\begin{grammarlens}[title={What you've built}]
One more word for saying what happens if.
\end{grammarlens}

\subsection*{Your turn}
\begin{itemize}
\raggedright
  \item \emph{Say it:} \emph{dāyaḥ}
  \item \emph{Say it:} \emph{dāyaḥ}, once more
  \item \emph{Say it:} \emph{yadi pituḥ dāyaḥ asti, tarhi bhrātaraḥ taṃ vaṇṭayanti.} — If there is an inheritance from their father, the brothers share it.
  \item \emph{Recall:} say the Sanskrit for divides up, then the Sanskrit for each one, then say \emph{dāyaḥ} again
\end{itemize}

\begin{tcolorbox}[breakable,colback=teal!4,colframe=teal!35!black,title={Before you move on}]
What does \sk{दायः} mean? (\textquotedblleft{}An inheritance\textquotedblright{}.) Say it once more.
\end{tcolorbox}
\section[\texorpdfstring{\sk{विनियुंक्ते} (viniyuṅkte)}{विनियुंक्ते (viniyuṅkte)}]{\sk{विनियुंक्ते} (viniyuṅkte) — puts to use}

\label{lesson:SA-C550-viniyunkte}



\begin{quote}
Before the new one: say the Sanskrit for each one, then the Sanskrit for an inheritance.
\end{quote}

\subsection*{You'll want to know: \sk{विनियुंक्ते}}
\textbf{\sk{विनियुंक्ते}} — \emph{viniyuṅkte} — \textquotedblleft{}puts to use\textquotedblright{}.

\begin{quote}
\sk{यदि} \sk{सः} \sk{धनं} \sk{सम्यक्} \sk{विनियुंक्ते}, \sk{तर्हि} \sk{सः} \sk{धनिकः} \sk{भविष्यति।}
\end{quote}

\begin{quote}
\emph{yadi saḥ dhanaṃ samyak viniyuṅkte, tarhi saḥ dhanikaḥ bhaviṣyati.}
\end{quote}

\begin{quote}
If he puts his money to good use, he will become rich.
\end{quote}

Past (he …): \sk{सः} \sk{धनं} \sk{सम्यक्} \sk{विनियुक्तवान्}, \emph{saḥ dhanaṃ samyak viniyuktavān}.

\begin{cousinweb}
\textbf{\sk{विनियुंक्ते}} (\emph{viniyuṅkte}) is built on \emph{vi}, \emph{ni} and the root √yuj, \textquotedblleft{}to yoke\textquotedblright{}; MW gives \emph{viniyuj} as to separate, and to assign, to commit, to appoint to a use (MW).
\end{cousinweb}

\begin{grammarlens}[title={What you've built}]
One more word for saying what happens if.
\end{grammarlens}

\subsection*{Your turn}
\begin{itemize}
\raggedright
  \item \emph{Say it:} \emph{viniyuṅkte}
  \item \emph{Say it:} \emph{viniyuṅkte}, once more
  \item \emph{Say it:} \emph{yadi saḥ dhanaṃ samyak viniyuṅkte, tarhi saḥ dhanikaḥ bhaviṣyati.} — If he puts his money to good use, he will become rich.
  \item \emph{Recall:} say the Sanskrit for each one, then the Sanskrit for an inheritance, then say \emph{viniyuṅkte} again
\end{itemize}

\begin{tcolorbox}[breakable,colback=teal!4,colframe=teal!35!black,title={Before you move on}]
What does \sk{विनियुंक्ते} mean? (\textquotedblleft{}Puts to use\textquotedblright{}.) Say it once more.
\end{tcolorbox}
\section[\texorpdfstring{\sk{वितरणम्} (vitaraṇam)}{वितरणम् (vitaraṇam)}]{\sk{वितरणम्} (vitaraṇam) — a handing out}

\label{lesson:SA-C550-vitaranam}



\begin{quote}
Before the new one: say the Sanskrit for an inheritance, then the Sanskrit for puts to use.
\end{quote}

\subsection*{You'll want to know: \sk{वितरणम्}}
\textbf{\sk{वितरणम्}} — \emph{vitaraṇam} — \textquotedblleft{}a handing out\textquotedblright{}.

\begin{quote}
\sk{यदि} \sk{उत्सवे} \sk{वितरणं} \sk{भवति}, \sk{तर्हि} \sk{बालाः} \sk{प्रसन्नाः} \sk{सन्ति।}
\end{quote}

\begin{quote}
\emph{yadi utsave vitaraṇaṃ bhavati, tarhi bālāḥ prasannāḥ santi.}
\end{quote}

\begin{quote}
If things are handed out at the festival, the children are happy.
\end{quote}

A neuter noun.

\begin{cousinweb}
\textbf{\sk{वितरणम्}} (\emph{vitaraṇam}) is built on \emph{vi} and the root √tṝ, \textquotedblleft{}to cross\textquotedblright{}, which also gives \emph{vitarati}, \textquotedblleft{}distributes\textquotedblright{}; MW gives \emph{vitaraṇa} as leading over, and granting, bestowal, a gift (MW).
\end{cousinweb}

\begin{grammarlens}[title={What you've built}]
One more word for saying what happens if.
\end{grammarlens}

\subsection*{Your turn}
\begin{itemize}
\raggedright
  \item \emph{Say it:} \emph{vitaraṇam}
  \item \emph{Say it:} \emph{vitaraṇam}, once more
  \item \emph{Say it:} \emph{yadi utsave vitaraṇaṃ bhavati, tarhi bālāḥ prasannāḥ santi.} — If things are handed out at the festival, the children are happy.
  \item \emph{Recall:} say the Sanskrit for an inheritance, then the Sanskrit for puts to use, then say \emph{vitaraṇam} again
\end{itemize}

\begin{tcolorbox}[breakable,colback=teal!4,colframe=teal!35!black,title={Before you move on}]
What does \sk{वितरणम्} mean? (\textquotedblleft{}A handing out\textquotedblright{}.) Say it once more.
\end{tcolorbox}
